\section{Scope, conventions, and the status of the results}
\label{sec:scope}

The starting point is the unified reading manuscript
\emph{Polylogarithms and their Arithmetic Bridges} \cite{manuscript}.
Its purpose is broader than the evaluation of isolated constants:
classical polylogarithms, colored nested sums, gamma functions,
Hurwitz-zeta jets, and algebraic ladders are different coordinates on
related analytic constructions. The present continuation selects three
places where those coordinates can be turned into short, checkable
proofs: the same-argument Gaussian depth problem, algebraic unit ladders,
and log-gamma moments.

The mathematical task is to explain why an identity holds. A small
decimal residual is useful for detecting a sign or transcription error,
but it cannot replace a derivation. Accordingly, every proved identity
below either follows from a displayed analytic argument or is accompanied
by exact rational operations on explicitly stated functional equations.
The remaining conjectures are kept separate from these results.

\subsection{What changes in the manuscript}

Table~\ref{tab:status} lists the precise changes of mathematical status.
The reference to a chapter is to the pinned source manuscript, not to
the section numbering of this article.

\begin{table}[htbp]
\centering\small
\begin{tabular}{>{\raggedright\arraybackslash}p{0.22\linewidth}>{\raggedright\arraybackslash}p{0.28\linewidth}>{\raggedright\arraybackslash}p{0.39\linewidth}}
\toprule
Source location & Previously recorded status & Result established here\\
\midrule
Chapter 4, weight-five Gaussian relation
& Numerical relation among four same-argument doubles
& Proof from two shuffle products; a spanning set of at most two quantities\\
\addlinespace
Chapter 4, five weight-six doubles
& Numerical evaluations
& All five proved by one explicit all-weight parity formula\\
\addlinespace
Chapter 4, three weight-four triples
& Numerical evaluations at \(i\)
& All three proved; an arbitrary-depth classical-reduction family\\
\addlinespace
Chapter 4, weight-six triple quotient
& Finite rank seven on ten coordinates
& Exact all-weight, all-depth formal theorem; the three stated coordinates are Lyndon generators\\
\addlinespace
Chapter 3, cubic and quartic bases
& Two trilogarithm formulas with numerical evidence
& Explicit functional-equation certificates for both\\
\addlinespace
Chapter 7, cubic log-gamma moment
& A named depth-two coordinate still sought
& Published Tornheim-derivative evaluation identified; certified independent numerical enclosure\\
\addlinespace
Log-gamma moment continuation
& Fixed-order factorial asymptotic expansion known in the literature
& Coefficient growth, divergence, and a remainder bound for a growing truncation order\\
\bottomrule
\end{tabular}
\caption{Exact scope of this continuation. None of the spanning statements
asserts numerical linear independence.}
\label{tab:status}
\end{table}

The eleven promoted identities are counted as one weight-five relation,
five weight-six relations, three triple relations, and two algebraic-base
relations. The extra weight-five shuffle equation is a further proved
identity and is what sharpens the spanning statement.
We do not count a repeated specialization of a general formula as a
separate solution of an open problem.

\subsection{Normalization and convergence}

We use decreasing summation indices:
\[
\Li_{s_1,\ldots,s_d}(z_1,\ldots,z_d)
=\sum_{n_1>\cdots>n_d>0}
\frac{z_1^{n_1}\cdots z_d^{n_d}}
{n_1^{s_1}\cdots n_d^{s_d}}.
\]
For a single varying argument, write
\[
F_{s_1,\ldots,s_d}(z)
=\Li_{s_1,\ldots,s_d}(z,1,\ldots,1).
\]
Where no confusion is possible we suppress the \(F\) and the trailing
ones. The weight is \(w=s_1+\cdots+s_d\), and the series depth is \(d\).
For \(|z|<1\), the series defining this one-variable family converges
absolutely. At a unit-circle point other than \(1\), the boundary values
used below agree with the ordinary convergent series; the needed
conditional double cases are justified in the parity proof. At higher
depth the inner multiple harmonic sum is \(O(\log^{d-1}n)\); its positive
successive increments, together with the factor \(n^{-s_1}\), give a
sequence tending to zero with finite total variation. Summation by parts
then proves convergence away from \(z=1\).

The forms
\[
\omega_0=\frac{\dd t}{t},\qquad
\omega_1=\frac{\dd t}{1-t}
\]
give the outer-letter-first word
\[
F_{s_1,\ldots,s_d}(z)
=I_{0^{s_1-1}1\cdots0^{s_d-1}1}(z).
\]
Here \(I_{a_1\cdots a_w}(z)=\int_0^z\omega_{a_1}(t)
I_{a_2\cdots a_w}(t)\), and \(I_{\varnothing}=1\).
Every displayed word ends in \(1\), so the origin is integrable.
For example, \(I_{01}=\Li_2\),
\(I_{011}=F_{2,1}\), and \(I_{101}=F_{1,2}\).
These examples fix the convention before any shuffle is used.

At the Gaussian point define
\[
g_{a,b}=\Im F_{a,b}(i),\quad
h_{a,b}=\Re F_{a,b}(i),\quad
G=\beta(2),\quad L=\log2,\quad
\lambda_n=\Im\Li_n((1+i)/2).
\]
The Dirichlet beta function is
\(\beta(s)=\sum_{n\ge0}(-1)^n/(2n+1)^s\).
All complex logarithms in the explicit half-point calculation are
principal values, with their continuation paths stated.
Logarithms of positive algebraic numbers in the ladder section are real.

\subsection{Prior results and the novelty boundary}

The general parity theorem is due to Panzer \cite{Panzer}, whose paper
already gives explicit formulas at depths two and three.
Umezawa \cite{Umezawa} subsequently supplied an explicit general-depth
formula, including regularized boundary formulations.
Our binomial expression is an independently proved specialization suited
to the manuscript's convention and to small exact certificates.
It is not a new general parity theorem.

The polynomial structure of a shuffle algebra on Lyndon generators is
classical \cite{Radford}; the rank formula below is its precise
application to the relevant one-variable subalgebra.
The Landen, Rogers, and Kummer equations used for the ladders are likewise
classical \cite{Zagier,NakamuraShiraishi}. What is supplied here is a
complete local route from those equations to the exact coefficients
recorded in the manuscript.

For log-gamma moments, the fixed-order factorial expansion was proved by
Amdeberhan et al.\ \cite{ACEMKM}, and the cubic Tornheim representation
was obtained by Bailey, Borwein, and Borwein \cite{BBB}.
The further coefficient asymptotic and the explicit growing-order
remainder estimate are developed here from those starting points.
These statements are presented with proofs and a limited novelty claim:
they extend the audited manuscript, while no assertion of exhaustive
priority over the entire inverse-gamma literature is made.
